\IfFileExists{stacks-project.cls}{%
\documentclass{stacks-project}
}{%
\documentclass{amsart}
}

% For dealing with references we use the comment environment
\usepackage{verbatim}
\newenvironment{reference}{\comment}{\endcomment}
%\newenvironment{reference}{}{}
\newenvironment{slogan}{\comment}{\endcomment}
\newenvironment{history}{\comment}{\endcomment}

% For commutative diagrams we use Xy-pic
\usepackage[all]{xy}

% We use 2cell for 2-commutative diagrams.
\xyoption{2cell}
\UseAllTwocells

% We use multicol for the list of chapters between chapters
\usepackage{multicol}

% This is generally recommended for better output
\usepackage{lmodern}
\usepackage[T1]{fontenc}

% For cross-file-references
\usepackage{xr-hyper}

% Package for hypertext links:
\usepackage{hyperref}

% For any local file, say "hello.tex" you want to link to please
% use \externaldocument[hello-]{hello}
\externaldocument[introduction-]{introduction}
\externaldocument[conventions-]{conventions}
\externaldocument[sets-]{sets}
\externaldocument[categories-]{categories}
\externaldocument[topology-]{topology}
\externaldocument[sheaves-]{sheaves}
\externaldocument[sites-]{sites}
\externaldocument[stacks-]{stacks}
\externaldocument[fields-]{fields}
\externaldocument[algebra-]{algebra}
\externaldocument[brauer-]{brauer}
\externaldocument[homology-]{homology}
\externaldocument[derived-]{derived}
\externaldocument[simplicial-]{simplicial}
\externaldocument[more-algebra-]{more-algebra}
\externaldocument[smoothing-]{smoothing}
\externaldocument[modules-]{modules}
\externaldocument[sites-modules-]{sites-modules}
\externaldocument[injectives-]{injectives}
\externaldocument[cohomology-]{cohomology}
\externaldocument[sites-cohomology-]{sites-cohomology}
\externaldocument[dga-]{dga}
\externaldocument[dpa-]{dpa}
\externaldocument[sdga-]{sdga}
\externaldocument[hypercovering-]{hypercovering}
\externaldocument[schemes-]{schemes}
\externaldocument[constructions-]{constructions}
\externaldocument[properties-]{properties}
\externaldocument[morphisms-]{morphisms}
\externaldocument[coherent-]{coherent}
\externaldocument[divisors-]{divisors}
\externaldocument[limits-]{limits}
\externaldocument[varieties-]{varieties}
\externaldocument[topologies-]{topologies}
\externaldocument[descent-]{descent}
\externaldocument[perfect-]{perfect}
\externaldocument[more-morphisms-]{more-morphisms}
\externaldocument[flat-]{flat}
\externaldocument[groupoids-]{groupoids}
\externaldocument[more-groupoids-]{more-groupoids}
\externaldocument[etale-]{etale}
\externaldocument[chow-]{chow}
\externaldocument[intersection-]{intersection}
\externaldocument[pic-]{pic}
\externaldocument[weil-]{weil}
\externaldocument[adequate-]{adequate}
\externaldocument[dualizing-]{dualizing}
\externaldocument[duality-]{duality}
\externaldocument[discriminant-]{discriminant}
\externaldocument[derham-]{derham}
\externaldocument[local-cohomology-]{local-cohomology}
\externaldocument[algebraization-]{algebraization}
\externaldocument[curves-]{curves}
\externaldocument[resolve-]{resolve}
\externaldocument[models-]{models}
\externaldocument[functors-]{functors}
\externaldocument[equiv-]{equiv}
\externaldocument[pione-]{pione}
\externaldocument[etale-cohomology-]{etale-cohomology}
\externaldocument[proetale-]{proetale}
\externaldocument[relative-cycles-]{relative-cycles}
\externaldocument[more-etale-]{more-etale}
\externaldocument[trace-]{trace}
\externaldocument[crystalline-]{crystalline}
\externaldocument[spaces-]{spaces}
\externaldocument[spaces-properties-]{spaces-properties}
\externaldocument[spaces-morphisms-]{spaces-morphisms}
\externaldocument[decent-spaces-]{decent-spaces}
\externaldocument[spaces-cohomology-]{spaces-cohomology}
\externaldocument[spaces-limits-]{spaces-limits}
\externaldocument[spaces-divisors-]{spaces-divisors}
\externaldocument[spaces-over-fields-]{spaces-over-fields}
\externaldocument[spaces-topologies-]{spaces-topologies}
\externaldocument[spaces-descent-]{spaces-descent}
\externaldocument[spaces-perfect-]{spaces-perfect}
\externaldocument[spaces-more-morphisms-]{spaces-more-morphisms}
\externaldocument[spaces-flat-]{spaces-flat}
\externaldocument[spaces-groupoids-]{spaces-groupoids}
\externaldocument[spaces-more-groupoids-]{spaces-more-groupoids}
\externaldocument[bootstrap-]{bootstrap}
\externaldocument[spaces-pushouts-]{spaces-pushouts}
\externaldocument[spaces-chow-]{spaces-chow}
\externaldocument[groupoids-quotients-]{groupoids-quotients}
\externaldocument[spaces-more-cohomology-]{spaces-more-cohomology}
\externaldocument[spaces-simplicial-]{spaces-simplicial}
\externaldocument[spaces-duality-]{spaces-duality}
\externaldocument[formal-spaces-]{formal-spaces}
\externaldocument[restricted-]{restricted}
\externaldocument[spaces-resolve-]{spaces-resolve}
\externaldocument[formal-defos-]{formal-defos}
\externaldocument[defos-]{defos}
\externaldocument[cotangent-]{cotangent}
\externaldocument[examples-defos-]{examples-defos}
\externaldocument[algebraic-]{algebraic}
\externaldocument[examples-stacks-]{examples-stacks}
\externaldocument[stacks-sheaves-]{stacks-sheaves}
\externaldocument[criteria-]{criteria}
\externaldocument[artin-]{artin}
\externaldocument[quot-]{quot}
\externaldocument[stacks-properties-]{stacks-properties}
\externaldocument[stacks-morphisms-]{stacks-morphisms}
\externaldocument[stacks-limits-]{stacks-limits}
\externaldocument[stacks-cohomology-]{stacks-cohomology}
\externaldocument[stacks-perfect-]{stacks-perfect}
\externaldocument[stacks-introduction-]{stacks-introduction}
\externaldocument[stacks-more-morphisms-]{stacks-more-morphisms}
\externaldocument[stacks-geometry-]{stacks-geometry}
\externaldocument[moduli-]{moduli}
\externaldocument[moduli-curves-]{moduli-curves}
\externaldocument[examples-]{examples}
\externaldocument[exercises-]{exercises}
\externaldocument[guide-]{guide}
\externaldocument[desirables-]{desirables}
\externaldocument[coding-]{coding}
\externaldocument[obsolete-]{obsolete}
\externaldocument[fdl-]{fdl}
\externaldocument[index-]{index}

% Theorem environments.
%
\theoremstyle{plain}
\newtheorem{theorem}[subsection]{Theorem}
\newtheorem{proposition}[subsection]{Proposition}
\newtheorem{lemma}[subsection]{Lemma}

\theoremstyle{definition}
\newtheorem{definition}[subsection]{Definition}
\newtheorem{example}[subsection]{Example}
\newtheorem{exercise}[subsection]{Exercise}
\newtheorem{situation}[subsection]{Situation}

\theoremstyle{remark}
\newtheorem{remark}[subsection]{Remark}
\newtheorem{remarks}[subsection]{Remarks}

\numberwithin{equation}{subsection}

% Macros
%
\def\lim{\mathop{\mathrm{lim}}\nolimits}
\def\colim{\mathop{\mathrm{colim}}\nolimits}
\def\Spec{\mathop{\mathrm{Spec}}}
\def\Hom{\mathop{\mathrm{Hom}}\nolimits}
\def\Ext{\mathop{\mathrm{Ext}}\nolimits}
\def\SheafHom{\mathop{\mathcal{H}\!\mathit{om}}\nolimits}
\def\SheafExt{\mathop{\mathcal{E}\!\mathit{xt}}\nolimits}
\def\Sch{\mathit{Sch}}
\def\Mor{\mathop{\mathrm{Mor}}\nolimits}
\def\Ob{\mathop{\mathrm{Ob}}\nolimits}
\def\Sh{\mathop{\mathit{Sh}}\nolimits}
\def\NL{\mathop{N\!L}\nolimits}
\def\CH{\mathop{\mathrm{CH}}\nolimits}
\def\proetale{{pro\text{-}\acute{e}tale}}
\def\etale{{\acute{e}tale}}
\def\QCoh{\mathit{QCoh}}
\def\Ker{\mathop{\mathrm{Ker}}}
\def\Im{\mathop{\mathrm{Im}}}
\def\Coker{\mathop{\mathrm{Coker}}}
\def\Coim{\mathop{\mathrm{Coim}}}

% Boxtimes
%
\DeclareMathSymbol{\boxtimes}{\mathbin}{AMSa}{"02}

%
% Macros for moduli stacks/spaces
%
\def\QCohstack{\mathcal{QC}\!\mathit{oh}}
\def\Cohstack{\mathcal{C}\!\mathit{oh}}
\def\Spacesstack{\mathcal{S}\!\mathit{paces}}
\def\Quotfunctor{\mathrm{Quot}}
\def\Hilbfunctor{\mathrm{Hilb}}
\def\Curvesstack{\mathcal{C}\!\mathit{urves}}
\def\Polarizedstack{\mathcal{P}\!\mathit{olarized}}
\def\Complexesstack{\mathcal{C}\!\mathit{omplexes}}
% \Pic is the operator that assigns to X its picard group, usage \Pic(X)
% \Picardstack_{X/B} denotes the Picard stack of X over B
% \Picardfunctor_{X/B} denotes the Picard functor of X over B
\def\Pic{\mathop{\mathrm{Pic}}\nolimits}
\def\Picardstack{\mathcal{P}\!\mathit{ic}}
\def\Picardfunctor{\mathrm{Pic}}
\def\Deformationcategory{\mathcal{D}\!\mathit{ef}}
\usepackage{cleveref}
\crefname{originallemma}{original lemma}{original lemmas}
\Crefname{originallemma}{Original lemma}{Original lemmas}
\makeatletter
\newlabel{categories-lemma-limits-products-equalizers}{{002N}{source}{}{sourcepointer-1}{}}
\newlabel{categories-lemma-limits-products-equalizers@cref}{{[originallemma][0][]002N}{[1][0][]source}{}{}{}}
\makeatother
\crefname{originallemma}{original lemma}{original lemmas}
\Crefname{originallemma}{Original lemma}{Original lemmas}
\makeatletter
\newlabel{homology-lemma-good-resolution-gives-qis}{{09IZ}{source}{}{sourcepointer-2}{}}
\newlabel{homology-lemma-good-resolution-gives-qis@cref}{{[originallemma][0][]09IZ}{[1][0][]source}{}{}{}}
\makeatother
\crefname{originallemma}{original lemma}{original lemmas}
\Crefname{originallemma}{Original lemma}{Original lemmas}
\makeatletter
\newlabel{homology-lemma-good-right-resolution-gives-qis}{{09J0}{source}{}{sourcepointer-3}{}}
\newlabel{homology-lemma-good-right-resolution-gives-qis@cref}{{[originallemma][0][]09J0}{[1][0][]source}{}{}{}}
\makeatother
\crefname{originallemma}{original lemma}{original lemmas}
\Crefname{originallemma}{Original lemma}{Original lemmas}
\makeatletter
\newlabel{derived-lemma-representable-homological}{{0149}{source}{}{sourcepointer-4}{}}
\newlabel{derived-lemma-representable-homological@cref}{{[originallemma][0][]0149}{[1][0][]source}{}{}{}}
\makeatother
\crefname{originallemma}{original lemma}{original lemmas}
\Crefname{originallemma}{Original lemma}{Original lemmas}
\makeatletter
\newlabel{derived-lemma-easier-axiom-four}{{014C}{source}{}{sourcepointer-5}{}}
\newlabel{derived-lemma-easier-axiom-four@cref}{{[originallemma][0][]014C}{[1][0][]source}{}{}{}}
\makeatother
\crefname{originallemma}{original lemma}{original lemmas}
\Crefname{originallemma}{Original lemma}{Original lemmas}
\makeatletter
\newlabel{derived-lemma-make-injective}{{013N}{source}{}{sourcepointer-6}{}}
\newlabel{derived-lemma-make-injective@cref}{{[originallemma][0][]013N}{[1][0][]source}{}{}{}}
\makeatother
\title{The derived category of a dg algebra has direct sums and products computed levelwise}
\author{The Stacks Project Authors}
\date{stacks-project:dga.tex:local-snapshot-20261012; GNU FDL 1.2 or later}
\begin{document}
\maketitle
\noindent\textit{Editorial note.} The counted statement is the author's own lemma environment labelled lemma-derived-products (source lines 2894-2904); the proof environment begins two lines below it and is its own proof, closing at source line 2950.\par\medskip
\noindent\textit{Editorial note.} The proof invokes the author's derived-Hom computation by label; that lemma is retained as a context span and the generator expands the context transitively (budget 400 lines per the Lead's ruling), while reference targets in other chapters that carry canonical project tags are auto-declared as non-load-bearing source pointers.\par\medskip
\noindent\textit{Editorial note.} The retained preamble is the chapter's own minimal preamble (sources/stacks-dga/preamble.tex, taken with its real line count), a verbatim excerpt of the native Xy-pic setup, the native theorem declarations with the mathematical macros, and the project's own no-op reference environments; no macro was invented or redefined.\par\medskip
\noindent\textit{Editorial scope.} One original result of the Stacks Project's Differential Graded Algebra chapter, counted once: the derived category of a dg algebra has direct sums and products, computed by the corresponding operations on dg modules.\par\medskip
\noindent\textit{Allowed standard prerequisites (editorial).}\begin{itemize}
\item Derived category of a dg algebra. The localisation of the homotopy category of dg modules at quasi-isomorphisms, with morphisms computed as a derived Hom.. It is the category in which direct sums and products are asserted to exist (source lines 2896-2904, 2906-2935).
\item Direct sums and products of dg modules. The degreewise direct sum and product of a family of dg modules, with the induced differentials, and their universal properties.. The candidates are the levelwise constructions and their universal properties have to be transferred (source lines 2906-2948).
\item Derived Hom and quasi-isomorphism invariance. Morphisms in the derived category are computed by the derived Hom, which only depends on quasi-isomorphism classes (the author's lemma lemma-hom-derived, retained as context).. It is what allows the universal property to be checked up to quasi-isomorphism (source lines 2910-2947).
\end{itemize}\medskip
\ifdefined\hypertarget\hypertarget{sourcepointer-1}{}\fi
\noindent\textit{Source reference.} Original lemma 002N: sources/stacks-chow/tags.txt, source lines 104--104. Cross-chapter pointer whose target carries the canonical tag 002N in the project tag list; declared so the reference resolves inside the assembled package, while every step the counted argument uses remains inside the retained spans.\par\medskip
\ifdefined\hypertarget\hypertarget{sourcepointer-2}{}\fi
\noindent\textit{Source reference.} Original lemma 09IZ: sources/stacks-chow/tags.txt, source lines 12057--12057. Cross-chapter pointer whose target carries the canonical tag 09IZ in the project tag list; declared so the reference resolves inside the assembled package, while every step the counted argument uses remains inside the retained spans.\par\medskip
\ifdefined\hypertarget\hypertarget{sourcepointer-3}{}\fi
\noindent\textit{Source reference.} Original lemma 09J0: sources/stacks-chow/tags.txt, source lines 12058--12058. Cross-chapter pointer whose target carries the canonical tag 09J0 in the project tag list; declared so the reference resolves inside the assembled package, while every step the counted argument uses remains inside the retained spans.\par\medskip
\ifdefined\hypertarget\hypertarget{sourcepointer-4}{}\fi
\noindent\textit{Source reference.} Original lemma 0149: sources/stacks-chow/tags.txt, source lines 1464--1464. Cross-chapter pointer whose target carries the canonical tag 0149 in the project tag list; declared so the reference resolves inside the assembled package, while every step the counted argument uses remains inside the retained spans.\par\medskip
\ifdefined\hypertarget\hypertarget{sourcepointer-5}{}\fi
\noindent\textit{Source reference.} Original lemma 014C: sources/stacks-chow/tags.txt, source lines 1467--1467. Cross-chapter pointer whose target carries the canonical tag 014C in the project tag list; declared so the reference resolves inside the assembled package, while every step the counted argument uses remains inside the retained spans.\par\medskip
\ifdefined\hypertarget\hypertarget{sourcepointer-6}{}\fi
\noindent\textit{Source reference.} Original lemma 013N: sources/stacks-chow/tags.txt, source lines 1442--1442. Cross-chapter pointer whose target carries the canonical tag 013N in the project tag list; declared so the reference resolves inside the assembled package, while every step the counted argument uses remains inside the retained spans.\par\medskip
\begin{lemma}
\label{lemma-hom-derived}
Let $(A, \text{d})$ be a differential graded algebra.
Let $M$ and $N$ be differential graded $A$-modules.
\begin{enumerate}
\item Let $P \to M$ be a P-resolution as in
Lemma \ref{lemma-resolve}. Then
$$
\Hom_{D(A, \text{d})}(M, N) =
\Hom_{K(\text{Mod}_{(A, \text{d})})}(P, N)
$$
\item Let $N \to I$ be an I-resolution as in
Lemma \ref{lemma-right-resolution}. Then
$$
\Hom_{D(A, \text{d})}(M, N) =
\Hom_{K(\text{Mod}_{(A, \text{d})})}(M, I)
$$
\end{enumerate}
\end{lemma}

\begin{proof}
Let $P \to M$ be as in (1). Since $P \to M$ is a quasi-isomorphism we see that
$$
\Hom_{D(A, \text{d})}(P, N) = \Hom_{D(A, \text{d})}(M, N)
$$
by definition of the derived category. A morphism
$f : P \to N$ in $D(A, \text{d})$ is equal to
$s^{-1}f'$ where $f' : P \to N'$ is a morphism and
$s : N \to N'$ is a quasi-isomorphism. Choose a distinguished triangle
$$
N \to N' \to Q \to N[1]
$$
As $s$ is a quasi-isomorphism, we see that $Q$ is acyclic. Thus
$\Hom_{K(\text{Mod}_{(A, \text{d})})}(P, Q[k]) = 0$ for all $k$ by
Lemma \ref{lemma-property-P-K-projective}. Since
$\Hom_{K(\text{Mod}_{(A, \text{d})})}(P, -)$
is cohomological, we conclude that we can lift $f' : P \to N'$
uniquely to a morphism $f : P \to N$. This finishes the proof.

\medskip\noindent
The proof of (2) is dual to that of (1) using
Lemma \ref{lemma-property-I-K-injective} in stead of
Lemma \ref{lemma-property-P-K-projective}.
\end{proof}

\begin{lemma}
\label{lemma-dgm-abelian}
Let $(A, d)$ be a differential graded algebra. The category
$\text{Mod}_{(A, \text{d})}$ is abelian and has arbitrary limits and colimits.
\end{lemma}

\begin{proof}
Kernels and cokernels commute with taking underlying $A$-modules.
Similarly for direct sums and colimits. In other words, these operations
in $\text{Mod}_{(A, \text{d})}$ commute with the forgetful functor to the
category of $A$-modules. This is not the case for products and limits.
Namely, if $N_i$, $i \in I$ is a family of
differential graded $A$-modules, then the product $\prod N_i$ in
$\text{Mod}_{(A, \text{d})}$ is given by setting $(\prod N_i)^n = \prod N_i^n$
and $\prod N_i = \bigoplus_n (\prod N_i)^n$. Thus we see that the product
does commute with the forgetful functor to the category of graded $A$-modules.
A category with products and equalizers has limits, see
Categories, Lemma \ref{categories-lemma-limits-products-equalizers}.
\end{proof}

\begin{lemma}
\label{lemma-homotopy-direct-sums}
Let $(A, \text{d})$ be a differential graded algebra.
The homotopy category $K(\text{Mod}_{(A, \text{d})})$
has direct sums and products.
\end{lemma}

\begin{proof}
Omitted. Hint: Just use the direct sums and products as in
Lemma \ref{lemma-dgm-abelian}. This works because we saw that
these functors commute with the forgetful functor to the category
of graded $A$-modules and because $\prod$ is an exact functor
on the category of families of abelian groups.
\end{proof}

\begin{lemma}
\label{lemma-resolve}
Let $(A, \text{d})$ be a differential graded algebra.
Let $M$ be a differential graded $A$-module. There exists a homomorphism
$P \to M$ of differential graded $A$-modules such that
\begin{enumerate}
\item $P \to M$ is a quasi-isomorphism, and
\item $P$ has property (P).
\end{enumerate}
\end{lemma}

\begin{proof}
Set $M = M_0$. We inductively choose short exact sequences
$$
0 \to M_{i + 1} \to P_i \to M_i \to 0
$$
where the maps $P_i \to M_i$ are chosen as in Lemma \ref{lemma-good-quotient}.
This gives a ``resolution''
$$
\ldots \to P_2 \xrightarrow{f_2} P_1 \xrightarrow{f_1} P_0 \to M \to 0
$$
Then we set
$$
P = \bigoplus\nolimits_{i \geq 0} P_i
$$
as an $A$-module with grading given by
$P^n = \bigoplus_{a + b = n} P_{-a}^b$ and
differential (as in the construction of the total complex associated
to a double complex) by
$$
\text{d}_P(x) = f_{-a}(x) + (-1)^a \text{d}_{P_{-a}}(x)
$$
for $x \in P_{-a}^b$. With these conventions $P$ is indeed a differential
graded $A$-module. Recalling that each $P_i$ has a two step filtration
$0 \to P_i' \to P_i \to P_i'' \to 0$ we set
$$
F_{2i}P = \bigoplus\nolimits_{i \geq j \geq 0} P_j
\subset
\bigoplus\nolimits_{i \geq 0} P_i = P
$$
and we add $P'_{i + 1}$ to $F_{2i}P$ to get $F_{2i + 1}$.
These are differential graded submodules and the successive quotients
are direct sums of shifts of $A$. By
Lemma \ref{lemma-target-graded-projective} we see that
the inclusions $F_iP \to F_{i + 1}P$ are admissible monomorphisms.
Finally, we have to show that the map $P \to M$ (given by the
augmentation $P_0 \to M$) is a quasi-isomorphism. This follows from
Homology, Lemma \ref{homology-lemma-good-resolution-gives-qis}.
\end{proof}

\begin{lemma}
\label{lemma-right-resolution}
Let $(A, \text{d})$ be a differential graded algebra.
Let $M$ be a differential graded $A$-module. There exists a homomorphism
$M \to I$ of differential graded $A$-modules such that
\begin{enumerate}
\item $M \to I$ is a quasi-isomorphism, and
\item $I$ has property (I).
\end{enumerate}
\end{lemma}

\begin{proof}
Set $M = M_0$. We inductively choose short exact sequences
$$
0 \to M_i \to I_i \to M_{i + 1} \to 0
$$
where the maps $M_i \to I_i$ are chosen as in Lemma \ref{lemma-good-sub}.
This gives a ``resolution''
$$
0 \to M \to I_0 \xrightarrow{f_0} I_1 \xrightarrow{f_1} I_1 \to \ldots
$$
Denote $I$ the differential graded $A$-module with graded parts
$$
I^n = \prod\nolimits_{i \geq 0} I^{n - i}_i
$$
and differential defined by
$$
\text{d}_I(x) = f_i(x) + (-1)^i \text{d}_{I_i}(x)
$$
for $x \in I_i^{n - i}$. With these conventions $I$ is indeed a differential
graded $A$-module. Recalling that each $I_i$ has a two step filtration
$0 \to I_i' \to I_i \to I_i'' \to 0$ we set
$$
F_{2i}I^n = \prod\nolimits_{j \geq i} I^{n - j}_j
\subset
\prod\nolimits_{i \geq 0} I^{n - i}_i = I^n
$$
and we add a factor $I'_{i + 1}$ to $F_{2i}I$ to get $F_{2i + 1}I$.
These are differential graded submodules and the successive quotients
are products of shifts of $A^\vee$. By
Lemma \ref{lemma-source-graded-injective} we see that
the inclusions $F_{i + 1}I \to F_iI$ are admissible monomorphisms.
Finally, we have to show that the map $M \to I$ (given by the
augmentation $M \to I_0$) is a quasi-isomorphism. This follows from
Homology, Lemma \ref{homology-lemma-good-right-resolution-gives-qis}.
\end{proof}

\begin{lemma}
\label{lemma-property-I-K-injective}
Let $(A, \text{d})$ be a differential graded algebra.
Let $I$ be a differential graded $A$-module with property (I).
Then
$$
\Hom_{K(\text{Mod}_{(A, \text{d})})}(N, I) = 0
$$
for all acyclic differential graded $A$-modules $N$.
\end{lemma}

\begin{proof}
We will use that $K(\text{Mod}_{(A, \text{d})})$ is a triangulated
category (Proposition \ref{proposition-homotopy-category-triangulated}).
Let $F_\bullet$ be a filtration on $I$ as in property (I).
The short exact sequence of Lemma \ref{lemma-property-I-sequence}
produces a distinguished triangle. Hence by
Derived Categories, Lemma \ref{derived-lemma-representable-homological}
it suffices to show that
$$
\Hom_{K(\text{Mod}_{(A, \text{d})})}(N, I/F_iI) = 0
$$
for all acyclic differential graded $A$-modules $N$ and all $i$.
Each of the differential graded modules $I/F_iI$ has a finite filtration
by admissible monomorphisms, whose graded pieces are
products of $A^\vee[k]$. Thus it suffices to prove that
$$
\Hom_{K(\text{Mod}_{(A, \text{d})})}(N, A^\vee[k]) = 0
$$
for all acyclic differential graded $A$-modules $N$ and all $k$.
This follows from Lemma \ref{lemma-hom-into-shift-dual-free}
and the fact that $(-)^\vee$ is an exact functor.
\end{proof}

\begin{lemma}
\label{lemma-property-P-K-projective}
Let $(A, \text{d})$ be a differential graded algebra.
Let $P$ be a differential graded $A$-module with property (P).
Then
$$
\Hom_{K(\text{Mod}_{(A, \text{d})})}(P, N) = 0
$$
for all acyclic differential graded $A$-modules $N$.
\end{lemma}

\begin{proof}
We will use that $K(\text{Mod}_{(A, \text{d})})$ is a triangulated
category (Proposition \ref{proposition-homotopy-category-triangulated}).
Let $F_\bullet$ be a filtration on $P$ as in property (P).
The short exact sequence of Lemma \ref{lemma-property-P-sequence}
produces a distinguished triangle. Hence by
Derived Categories, Lemma \ref{derived-lemma-representable-homological}
it suffices to show that
$$
\Hom_{K(\text{Mod}_{(A, \text{d})})}(F_iP, N) = 0
$$
for all acyclic differential graded $A$-modules $N$ and all $i$.
Each of the differential graded modules $F_iP$ has a finite filtration
by admissible monomorphisms, whose graded pieces are direct sums
of shifts $A[k]$. Thus it suffices to prove that
$$
\Hom_{K(\text{Mod}_{(A, \text{d})})}(A[k], N) = 0
$$
for all acyclic differential graded $A$-modules $N$ and all $k$.
This follows from Lemma \ref{lemma-hom-from-shift-free}.
\end{proof}

\begin{lemma}
\label{lemma-good-quotient}
Let $(A, \text{d})$ be a differential graded algebra.
Let $M$ be a differential graded $A$-module. There exists a homomorphism
$P \to M$ of differential graded $A$-modules with the following
properties
\begin{enumerate}
\item $P \to M$ is surjective,
\item $\Ker(\text{d}_P) \to \Ker(\text{d}_M)$ is surjective, and
\item $P$ sits in an admissible short exact sequence
$0 \to P' \to P \to P'' \to 0$ where $P'$, $P''$ are direct sums
of shifts of $A$.
\end{enumerate}
\end{lemma}

\begin{proof}
Let $P_k$ be the free $A$-module with generators $x, y$ in degrees
$k$ and $k + 1$. Define the structure of a differential graded
$A$-module on $P_k$ by setting $\text{d}(x) = y$ and $\text{d}(y) = 0$.
For every element $m \in M^k$ there is a homomorphism
$P_k \to M$ sending $x$ to $m$ and $y$ to $\text{d}(m)$.
Thus we see that there is a surjection from a direct sum
of copies of $P_k$ to $M$. This clearly produces $P \to M$
having properties (1) and (3). To obtain property (2) note
that if $m \in \Ker(\text{d}_M)$ has degree $k$, then there is a map
$A[k] \to M$ mapping $1$ to $m$. Hence we can achieve (2) by adding
a direct sum of copies of shifts of $A$.
\end{proof}

\begin{lemma}
\label{lemma-target-graded-projective}
Let $(A, \text{d})$ be a differential graded algebra.
Let $M \to P$ be a surjective homomorphism of differential graded
$A$-modules. If $P$ is projective as a graded $A$-module, then
$M \to P$ is an admissible epimorphism.
\end{lemma}

\begin{proof}
This is immediate from the definitions.
\end{proof}

\begin{lemma}
\label{lemma-good-sub}
Let $(A, \text{d})$ be a differential graded algebra.
Let $M$ be a differential graded $A$-module. There exists a homomorphism
$M \to I$ of differential graded $A$-modules with the following
properties
\begin{enumerate}
\item $M \to I$ is injective,
\item $\Coker(\text{d}_M) \to \Coker(\text{d}_I)$ is injective,
and
\item $I$ sits in an admissible short exact sequence
$0 \to I' \to I \to I'' \to 0$ where $I'$, $I''$ are products
of shifts of $A^\vee$.
\end{enumerate}
\end{lemma}

\begin{proof}
We will use the functors $N \mapsto N^\vee$ (from left to right
differential graded modules and from right to left differential
graded modules) constructed in
Section \ref{section-modules-noncommutative-differential-graded}
and all of their properties.
For every $k \in \mathbf{Z}$ let $Q_k$ be the free left $A$-module with
generators $x, y$ in degrees $k$ and $k + 1$. Define the structure of a
left differential graded $A$-module on $Q_k$ by setting $\text{d}(x) = y$
and $\text{d}(y) = 0$. Arguing exactly as in the proof of
Lemma \ref{lemma-good-quotient} we find a surjection
$$
\bigoplus\nolimits_{i \in I} Q_{k_i} \longrightarrow M^\vee
$$
of left differential graded $A$-modules. Then we can consider the injection
$$
M \to (M^\vee)^\vee \to (\bigoplus\nolimits_{i \in I} Q_{k_i})^\vee =
\prod\nolimits_{i \in I} I_{k_i}
$$
where $I_k = Q_{-k}^\vee$ is the ``dual'' right differential graded $A$-module.
Further, the short exact sequence $0 \to A[-k - 1] \to Q_k \to A[-k] \to 0$
produces a short exact sequence
$0 \to A^\vee[k] \to I_k \to A^\vee[k + 1] \to 0$.

\medskip\noindent
The result of the previous paragraph produces $M \to I$
having properties (1) and (3). To obtain property (2), suppose
$\overline{m} \in \Coker(\text{d}_M)$ is a nonzero element of
degree $k$. Pick a map $\lambda : M^k \to \mathbf{Q}/\mathbf{Z}$
which vanishes on $\Im(M^{k - 1} \to M^k)$ but not on $m$. By
Lemma \ref{lemma-hom-into-shift-dual-free}
this corresponds to a homomorphism $M \to A^\vee[k]$ of
differential graded $A$-modules which does not vanish on $m$.
Hence we can achieve (2) by adding
a product of copies of shifts of $A^\vee$.
\end{proof}

\begin{lemma}
\label{lemma-source-graded-injective}
Let $(A, \text{d})$ be a differential graded algebra.
Let $I \to M$ be an injective homomorphism of differential graded
$A$-modules. If $I$ is graded injective, then
$I \to M$ is an admissible monomorphism.
\end{lemma}

\begin{proof}
This is immediate from the definitions.
\end{proof}

\begin{lemma}
\label{lemma-hom-into-shift-dual-free}
Let $(A, \text{d})$ be a differential graded algebra. Then we have
$$
\Hom_{\text{Mod}_{(A, \text{d})}}(M, A^\vee[k]) =
\Ker(\text{d} : (M^\vee)^k \to (M^\vee)^{k + 1})
$$
and
$$
\Hom_{K(\text{Mod}_{(A, \text{d})})}(M, A^\vee[k]) = H^k(M^\vee)
$$
as functors in the differential graded $A$-module $M$.
\end{lemma}

\begin{proof}
This is clear from the discussion above.
\end{proof}

\begin{lemma}
\label{lemma-property-I-sequence}
Let $(A, \text{d})$ be a differential graded algebra.
Let $I$ be a differential graded $A$-module. If $F_\bullet$
is a filtration as in property (I), then we obtain an
admissible short exact sequence
$$
0 \to I \to
\prod\nolimits I/F_iI \to
\prod\nolimits I/F_iI \to 0
$$
of differential graded $A$-modules.
\end{lemma}

\begin{proof}
Omitted. Hint: This is dual to Lemma \ref{lemma-property-P-sequence}.
\end{proof}

\begin{proposition}
\label{proposition-homotopy-category-triangulated}
Let $(A, \text{d})$ be a differential graded algebra. The homotopy category
$K(\text{Mod}_{(A, \text{d})})$ of differential graded $A$-modules with its
natural translation functors and distinguished triangles is a triangulated
category.
\end{proposition}

\begin{proof}
We know that $K(\text{Mod}_{(A, \text{d})})$ is a pre-triangulated category.
Hence it suffices to prove TR4 and to prove it we can use
Derived Categories, Lemma \ref{derived-lemma-easier-axiom-four}.
Let $K \to L$ and $L \to M$ be composable morphisms of
$K(\text{Mod}_{(A, \text{d})})$. By
Lemma \ref{lemma-sequence-maps-split} we may assume that
$K \to L$ and $L \to M$ are admissible monomorphisms.
In this case the result follows from
Lemma \ref{lemma-two-split-injections}.
\end{proof}

\begin{lemma}
\label{lemma-hom-from-shift-free}
Let $(A, d)$ be a differential graded algebra. Then we have
$$
\Hom_{\text{Mod}_{(A, \text{d})}}(A[k], M) =
\Ker(\text{d} : M^{-k} \to M^{-k + 1})
$$
and
$$
\Hom_{K(\text{Mod}_{(A, \text{d})})}(A[k], M) = H^{-k}(M)
$$
for any differential graded $A$-module $M$.
\end{lemma}

\begin{proof}
Immediate from the definitions.
\end{proof}

\begin{lemma}
\label{lemma-property-P-sequence}
Let $(A, \text{d})$ be a differential graded algebra.
Let $P$ be a differential graded $A$-module. If $F_\bullet$
is a filtration as in property (P), then we obtain an
admissible short exact sequence
$$
0 \to
\bigoplus\nolimits F_iP \to
\bigoplus\nolimits F_iP \to P \to 0
$$
of differential graded $A$-modules.
\end{lemma}

\begin{proof}
The second map is the direct sum of the inclusion maps.
The first map on the summand $F_iP$ of the source is the sum
of the identity $F_iP \to F_iP$ and the negative of the inclusion
map $F_iP \to F_{i + 1}P$. Choose homomorphisms $s_i : F_{i + 1}P \to F_iP$
of graded $A$-modules which are left inverse to the inclusion maps.
Composing gives maps $s_{j, i} : F_jP \to F_iP$ for all $j > i$.
Then a left inverse of the first arrow maps $x \in F_jP$ to
$(s_{j, 0}(x), s_{j, 1}(x), \ldots, s_{j, j - 1}(x), 0, \ldots)$
in $\bigoplus F_iP$.
\end{proof}

\section{Injective modules and differential graded algebras}
\label{section-modules-noncommutative-differential-graded}

\begin{lemma}
\label{lemma-sequence-maps-split}
Let $(A, \text{d})$ be a differential graded algebra.
Let $L_1 \to L_2 \to \ldots \to L_n$
be a sequence of composable homomorphisms of
differential graded $A$-modules.
There exists a commutative diagram
$$
\xymatrix{
L_1 \ar[r] &
L_2 \ar[r] &
\ldots \ar[r] &
L_n \\
M_1 \ar[r] \ar[u] &
M_2 \ar[r] \ar[u] &
\ldots \ar[r] &
M_n \ar[u]
}
$$
in $\text{Mod}_{(A, \text{d})}$ such that each $M_i \to M_{i + 1}$
is an admissible monomorphism and each $M_i \to L_i$
is a homotopy equivalence.
\end{lemma}

\begin{proof}
The case $n = 1$ is without content.
Lemma \ref{lemma-make-injective} is the case $n = 2$.
Suppose we have constructed the diagram
except for $M_n$. Apply Lemma \ref{lemma-make-injective} to
the composition $M_{n - 1} \to L_{n - 1} \to L_n$.
The result is a factorization $M_{n - 1} \to M_n \to L_n$
as desired.
\end{proof}

\begin{lemma}
\label{lemma-two-split-injections}
Let $(A, \text{d})$ be a differential graded algebra. Suppose that
$\alpha : K \to L$ and $\beta : L \to M$ are admissible monomorphisms
of differential graded $A$-modules. Then there exist distinguished triangles
$(K, L, Q_1, \alpha, p_1, d_1)$, $(K, M, Q_2, \beta \circ \alpha, p_2, d_2)$
and $(L, M, Q_3, \beta, p_3, d_3)$ for which TR4 holds.
\end{lemma}

\begin{proof}
Say $\pi_1 : L \to K$ and $\pi_3 : M \to L$ are homomorphisms
of graded $A$-modules which are left inverse to $\alpha$ and $\beta$.
Then also $K \to M$ is an admissible monomorphism with left
inverse $\pi_2 = \pi_1 \circ \pi_3$. Let us write $Q_1$, $Q_2$
and $Q_3$ for the cokernels of $K \to L$, $K \to M$, and $L \to M$.
Then we obtain identifications (as graded $A$-modules)
$Q_1 = \Ker(\pi_1)$, $Q_3 = \Ker(\pi_3)$ and
$Q_2 = \Ker(\pi_2)$. Then $L = K \oplus Q_1$ and
$M = L \oplus Q_3$ as graded $A$-modules. This implies
$M = K \oplus Q_1 \oplus Q_3$. Note that $\pi_2 = \pi_1 \circ \pi_3$
is zero on both $Q_1$ and $Q_3$. Hence $Q_2 = Q_1 \oplus Q_3$.
Consider the commutative diagram
$$
\begin{matrix}
0 & \to & K & \to & L & \to & Q_1 & \to & 0 \\
  &     & \downarrow &     & \downarrow &     & \downarrow  & \\
0 & \to & K & \to & M & \to & Q_2 & \to & 0 \\
  &     & \downarrow &     & \downarrow &     & \downarrow  & \\
0 & \to & L & \to & M & \to & Q_3 & \to & 0
\end{matrix}
$$
The rows of this diagram are admissible short exact sequences, and
hence determine distinguished triangles by definition. Moreover
downward arrows in the diagram above are compatible with the chosen
splittings and hence define morphisms of triangles
$$
(K \to L \to Q_1 \to K[1])
\longrightarrow
(K \to M \to Q_2 \to K[1])
$$
and
$$
(K \to M \to Q_2 \to K[1])
\longrightarrow
(L \to M \to Q_3 \to L[1]).
$$
Note that the splittings $Q_3 \to M$ of the bottom sequence in the
diagram provides a splitting for the split sequence
$0 \to Q_1 \to Q_2 \to Q_3 \to 0$ upon composing with $M \to Q_2$.
It follows easily from this that the morphism $\delta : Q_3 \to Q_1[1]$
in the corresponding distinguished triangle
$$
(Q_1 \to Q_2 \to Q_3 \to Q_1[1])
$$
is equal to the composition $Q_3 \to L[1] \to Q_1[1]$.
Hence we get a structure as in the conclusion of axiom TR4.
\end{proof}

\begin{lemma}
\label{lemma-make-injective}
Let $(A, \text{d})$ be a differential graded algebra.
Let $\alpha : K \to L$ be a homomorphism of differential graded
$A$-modules. There exists a factorization
$$
\xymatrix{
K \ar[r]^{\tilde \alpha} \ar@/_1pc/[rr]_\alpha &
\tilde L \ar[r]^\pi & L
}
$$
in $\text{Mod}_{(A, \text{d})}$ such that
\begin{enumerate}
\item $\tilde \alpha$ is an admissible monomorphism (see
Definition \ref{definition-admissible-ses}),
\item there is a morphism $s : L \to \tilde L$
such that $\pi \circ s = \text{id}_L$ and such that
$s \circ \pi$ is homotopic to $\text{id}_{\tilde L}$.
\end{enumerate}
\end{lemma}

\begin{proof}
The proof is identical to the proof of
Derived Categories, Lemma \ref{derived-lemma-make-injective}.
Namely, we set $\tilde L = L \oplus C(1_K)$ and we use elementary
properties of the cone construction.
\end{proof}

\begin{definition}
\label{definition-admissible-ses}
Let $(A, \text{d})$ be a differential graded algebra.
\begin{enumerate}
\item A homomorphism $K \to L$ of differential graded $A$-modules
is an {\it admissible monomorphism} if there exists a graded $A$-module
map $L \to K$ which is left inverse to $K \to L$.
\item A homomorphism $L \to M$ of differential graded $A$-modules
is an {\it admissible epimorphism} if there exists a graded $A$-module
map $M \to L$ which is right inverse to $L \to M$.
\item A short exact sequence $0 \to K \to L \to M \to 0$ of differential
graded $A$-modules is an {\it admissible short exact sequence}
if it is split as a sequence of graded $A$-modules.
\end{enumerate}
\end{definition}

\begin{lemma}
\label{lemma-derived-products}
Let $(A, \text{d})$ be a differential graded algebra. Then
\begin{enumerate}
\item $D(A, \text{d})$ has both direct sums and products,
\item direct sums are obtained by taking direct sums of differential graded
modules,
\item products are obtained by taking products of differential
graded modules.
\end{enumerate}
\end{lemma}

\begin{proof}
We will use that $\text{Mod}_{(A, \text{d})}$ is an abelian category
with arbitrary direct sums and products, and that these give rise
to direct sums and products in $K(\text{Mod}_{(A, \text{d})})$.
See Lemmas \ref{lemma-dgm-abelian} and \ref{lemma-homotopy-direct-sums}.

\medskip\noindent
Let $M_j$ be a family of differential graded $A$-modules.
Consider the graded direct sum $M = \bigoplus M_j$ which is a
differential graded $A$-module with the obvious.
For a differential graded $A$-module $N$ choose a quasi-isomorphism
$N \to I$ where $I$ is a differential graded $A$-module with property (I).
See Lemma \ref{lemma-right-resolution}.
Using Lemma \ref{lemma-hom-derived} we have
\begin{align*}
\Hom_{D(A, \text{d})}(M, N)
& =
\Hom_{K(A, \text{d})}(M, I) \\
& =
\prod \Hom_{K(A, \text{d})}(M_j, I) \\
& =
\prod \Hom_{D(A, \text{d})}(M_j, N)
\end{align*}
whence the existence of direct sums in $D(A, \text{d})$ as given in
part (2) of the lemma.

\medskip\noindent
Let $M_j$ be a family of differential graded $A$-modules.
Consider the product $M = \prod M_j$ of differential graded $A$-modules.
For a differential graded $A$-module $N$ choose a quasi-isomorphism
$P \to N$ where $P$ is a differential graded $A$-module with property (P).
See Lemma \ref{lemma-resolve}.
Using Lemma \ref{lemma-hom-derived} we have
\begin{align*}
\Hom_{D(A, \text{d})}(N, M)
& =
\Hom_{K(A, \text{d})}(P, M) \\
& =
\prod \Hom_{K(A, \text{d})}(P, M_j) \\
& =
\prod \Hom_{D(A, \text{d})}(N, M_j)
\end{align*}
whence the existence of direct sums in $D(A, \text{d})$ as given in
part (3) of the lemma.
\end{proof}

\end{document}
